\chapter{Introdução}
\label{cap:introducao}

% TODO: escrever texto de abertura do capitulo (1-2 paragrafos em alto nivel sobre o tema). Apresenta a fabrica de fios esmaltados e situa o leitor no problema que sera detalhado adiante. Conteudo previo desta secao esta disponivel no historico do git (commit anterior a af1cac4).


\section{Problema de pesquisa}
\label{sec:problema_pesquisa}

% TODO: enunciar o problema de pesquisa de forma autocontida.
% Cobertura prevista:
% - Mecanismos sobrepostos (materia-prima, equipamento, combinacao) que produzem fenotipos semelhantes na fratura, referenciando a Secao~\ref{sec:problema_rompimentos} do Cap. 2.
% - Ausencia de rotulo de causa nos dados operacionais, referenciando a Secao~\ref{sec:dados_disponiveis} do Cap. 2.
% - Enunciar a pergunta central:
%   "Como combinar sinais operacionais de maquina e evidencias estatisticas historicas de rompimentos para inferir a provavel origem de eventos de rompimento em uma fabrica de fios?"


\section{Objetivos}
\label{sec:objetivos}


\subsection{Objetivo geral}
\label{subsec:objetivo_geral}

% TODO (versao planejada): Desenvolver uma ferramenta hibrida de apoio a decisao que combine sinais operacionais de maquina e evidencias estatisticas historicas de rompimentos para inferir a provavel origem de eventos de rompimento em uma fabrica de fios esmaltados de cobre e aluminio.


\subsection{Objetivos específicos}
\label{subsec:objetivos_especificos}

% TODO (versao planejada com 5 itens):
% 1. Conduzir revisao sistematica da literatura sobre deteccao de anomalias e analise estatistica aplicadas a atribuicao de causa em processos industriais com ausencia de rotulos.
% 2. Caracterizar e preparar as fontes de dados operacionais e historicos disponiveis para analise.
% 3. Desenvolver um eixo de deteccao de anomalias operacionais a partir de series temporais de variaveis sensoriais.
% 4. Desenvolver um eixo de analise estatistica historica de rompimentos segmentado por lote, material, fornecedor e maquina.
% 5. Integrar os dois eixos em uma ferramenta de inferencia e avaliar o resultado a partir de casos historicos da fabrica.


\section{Estrutura do trabalho}
\label{sec:estrutura_trabalho}

% TODO: paragrafo descrevendo o que cada capitulo seguinte trata. Preencher depois que a estrutura final dos capitulos estiver consolidada (Cap. 2 Fundamentacao Teorica, Cap. 3 Metodologia, Cap. 4 Desenvolvimento da Ferramenta, Cap. 5 Resultados, Cap. 6 Conclusao).
